\chapter{Description du champ de vitesse}

\section{Description Eulérienne}

\subsection{Hypothèse de l'écoulement laminaire}
Nous allons un instant laisser en suspens les développements de l'équation du mouvement pour aborder la description du champ de vitesse.
Considérons un champ de vitesse qui, en toute généralité, dépend du temps : $\vec{u}(\vec{x}, t)$.
Tout d'abord, précisons que la notion même d'un champ de vitesse implique que nous pouvons définir, sans ambiguïté, une et une seule vitesse pour une position donnée, à tout moment. Or, nous avons vu que cette approche suppose d'une part une hypothèse de milieu continu, et l'existence de particules matérielles qui sont suffisamment à l'équilibre avec leur environnement pour être considérées comme homogènes. En particulier, nous sommes amenés à appliquer des opérateurs différentiels au champ de vitesse, ce qui exige que le champ de vitesse soit continûment dérivable.
Ces hypothèses nous amènent à considérer un écoulement suffisamment régulier que l'on qualifie d'\cindex{écoulement laminaire}, par opposition à l'écoulement turbulent. Les notions que nous allons développer ici supposent un écoulement laminaire.

\subsection{Lignes de courant}
À chaque temps $t$, nous pouvons définir un ensemble de courbes paramétriques $\vec{x}_{\vec{x}_0}$ définies par
\begin{equation}
  \od{\vec{x}}{s} = \vec{u}(\vec{x},t)\text{, avec $\vec{x}(0)=\vec{x}_0$}.
  \label{eq:ligne_courant}
\end{equation}
Il faut bien comprendre cette écriture car elle est piégeuse. Le paramètre $s$ est celui qui permet de tracer toute une courbe partant de $\vec{x}_0$, à un instant $t$ donné. La courbe est donc tracée en parcourant les valeurs de $s$ sur la droite réelle, en utilisant le champ $\vec{u}(\vec{x},t)$ à cet instant $t$ déterminé.

Supposons, pour fixer les idées, que le champ $\vec{u}$ soit invariant au cours du temps. Il s'agit donc d'un \cindex{écoulement stationnaire}. Réécrivons dans ce cas particulier :
\begin{displaymath}
  \od{\vec{x}}{s} = \vec{u}(\vec{x})\text{, avec $\vec{x}(0)=\vec{x}_0$}.
\end{displaymath}
Si on remplace, pour s'aider, le symbole $s$ par le symbole $t$, on voit immédiatement que la courbe $\vec{x}(s)$ répond à la même équation que celle qui détermine le chemin d'une particule de fluide au cours du temps, avec condition initiale $\vec{x}_0$ au temps $t=0$.
La courbe paramétrique définie par~\eqref{eq:ligne_courant} porte le nom de \cindex{ligne de courant}. Dans le cas particulier d'un champ de vitesse stationnaire, elle se superpose à la trajectoire ou chemin de toute particule passant à un moment donné par $\vec{x}_0$. Nous reviendrons sur cette distinction, mais soulignons dès à présent que la ligne de courant est toujours définie, même en régime non stationnaire.

\begin{key}{Les lignes de courant sont tangentes à la vitesse en tout point}
Exprimons l'équation~\eqref{eq:ligne_courant} selon ses composantes cartésiennes :
\begin{align*}
  \od{{x}}{s} &= u(\vec{x},t) &
  \od{{y}}{s} &= v(\vec{x},t) &
  \od{{z}}{s} &= w(\vec{x},t)
\end{align*}
On peut exprimer les incréments différentiels :
\begin{align*}
  \dif{x} &= u(\vec{x},t) \dif s &
  \dif{y} &= v(\vec{x},t) \dif s &
  \dif{z} &= w(\vec{x},t) \dif s.
\end{align*}
Ou encore :
\begin{align*}
  \frac{\dif{x}}{\dif{y}}  &= \frac{u(\vec{x},t)}{v(\vec{x},t)} &
  \frac{\dif{x}}{\dif{z}}  &= \frac{u(\vec{x},t)}{w(\vec{x},t)} &
  \frac{\dif{y}}{\dif{z}}  &= \frac{v(\vec{x},t)}{w(\vec{x},t)}
\end{align*}
ce qui revient simplement à écrire que la courbe décrite par $\vec{x}(s)$ est tangente, en tout point, au vecteur vitesse.
\end{key}

\begin{key}{Les lignes de courant ne peuvent se croiser que si la vitesse y est nulle}
Par l'absurde, cela impliquerait en effet qu'en un point et à un instant $t$, la vitesse peut prendre deux directions différentes, alors que l'hypothèse du milieu continu nous impose une vitesse bien définie en tout point.
\end{key}

Le cas de vitesse nulle est particulier : une particule de fluide n'atteindra jamais le point de vitesse nulle dans un écoulement stationnaire. Ces points sont des \cindex{points de stagnation}.

\subsection{Section de lignes de courant dans un écoulement}
Pour un champ de vitesse bidimensionnel $\vec{u}(x,y)$, considérons un ensemble de conditions initiales ${S}_{ab}(0)$ reliant deux points $a$ et $b$ de façon continue, et tel qu’il intercepte les lignes de courant de manière à ce que les vitesses soient toutes dans la même direction (par exemple, de gauche à droite).
On dit que ${S}_{ab}(0)$ constitue une \emph{section}, et plus précisément une \cindex{section connexe}, car elle ne contient pas de trou.
En construisant l’ensemble des lignes de courant associées à chacune de ces conditions initiales, on obtient une figure géométrique généralement bidimensionnelle. Pour tout paramètre $s$ (le temps restant figé), les ensembles ${S}_{ab}(s)$ constituent des sections correspondant à la projection continue des conditions initiales en fonction de $s$. Cette surface ainsi construite ne contiendra pas de trou non plus.
Pour le voir, introduisons la notion de \cindex{fonction de flot}, notée $\phi^s$, qui renvoie l’advection de la section initiale selon les lignes de courant :
\begin{equation}
  \phi^s({S}_{ab}(0)) = {S}_{ab}(s)
  \label{eq:flot_2d}
\end{equation}

Par conséquent, l’advection d’une section connexe par un flot continu reste connexe : \textbf{aucun trou n’est créé}.

\subsection{Écoulement tri-dimensionnel : tubes de courant}
Pour un champ de vitesse tridimensionnel $\vec{u}(x,y,z)$, nous pouvons définir un anneau ${C}_{ab}(0)$, c’est-à-dire une courbe fermée, dite \emph{simplement connexe}. L’advection de cette courbe selon le flot $\phi^s$ définit alors un \cindex{tube de courant} (Figure~\ref{fig:02_tube_courant}), dont la paroi est continue.

On appelle \emph{tube de courant} la surface formée par l’ensemble des lignes de courant issues d’une courbe fermée ${\partial S}(0)$ dans un champ de vitesse $\vec{u}(x,y,z)$. Autrement dit, si chaque point de $\partial S(0)$ est advecté selon le flot $\phi^s$, la surface latérale engendrée par l’ensemble des trajectoires forme un tube dont les lignes de courant constituent la paroi.

Les propriétés topologiques du flot garantissent la continuité de la paroi du tube : aucune ligne de courant ne traverse cette paroi, et donc aucun échange de matière ne se produit entre l’intérieur et l’extérieur du tube. De manière informelle, aucune ligne de courant ne peut s’échapper du tube : toutes restent confinées à l’intérieur de sa paroi.

\begin{key}{Conservation du débit massique en régime stationnaire}
En régime stationnaire, la masse $M$ du fluide contenu dans un tube de courant est constante au cours du temps. Autrement dit, la dérivée temporelle de la masse du tube est nulle :
\[
\frac{\partial M}{\partial t} = 0.
\]
Considérons deux sections transversales $S_1$ et $S_2$ du tube de courant, de normales respectives $\vec{n}_1$ et $\vec{n}_2$ orientées vers l’extérieur du volume délimité par ces sections (voir Figure~\ref{fig:02_tube_courant}). Le bilan de masse entre ces deux sections s’écrit :
\[
\int_{{S}_1} \rho \vec{u} \cdot (-\vec{n}_1) \, \dif S + \int_{{S}_2} \rho \vec{u} \cdot \vec{n}_2 \, \dif S = 0,
\]
où $\rho$ est la masse volumique du fluide. Les signes des normales reflètent le sens conventionnel du débit (entrant par ${S}_1$ et sortant par ${S}_2$).

En notant $Q_m = \int_{{S}} \rho \vec{u} \cdot \vec{n} \, \dif S$ le \cindex{débit massique} à travers une section ${S}$, cette équation se simplifie en :
\[
Q_{m,1} = Q_{m,2}.
\]
Ainsi, le débit massique est constant le long du tube de courant en régime stationnaire. Cette propriété est fondamentale pour l’étude des écoulements, et plus particulièrement les écoulements incompressibles car il permet, par exemple, de relier les vitesses moyennes dans des sections de tailles différentes.

  \begin{center}
   \captionsetup{type=figure}
    \input{Pdftex/02_tube_courant.pdf_tex}
  \caption{Le tube de courant est défini comme une surface imaginaire obtenue par l'ensemble continu de lignes de courant générées à partir d'une courbe fermée ${S}_{ab}(0)$. Les sections ${S}_1$ et ${S}_2$ illustrent les surfaces utilisées pour le bilan de masse.}\label{fig:02_tube_courant}
\end{center}


\end{key}

\subsection{Écoulements bi-dimensionnels}
Bien que nous vivions dans un monde tri-dimensionnel, beaucoup d'écoulements se décrivent très bien dans un espace à deux dimensions : les éléments de l'écoulement qui nous intéressent se trouvent dans un plan.
Ce sera typiquement le cas lorsque l'écoulement est confiné dans un domaine qui présente des dimensions très différentes selon l'horizontale et la verticale. C'est le cas des mers, des océans, ou dans certaines expériences de laboratoire comme celle de \cindex{Rayleigh-Bénard}\footnote{\url{https://www.youtube.com/watch?v=Eud7uG5JHng}}. De tels écoulements sont dits à coefficient d'aspect élevé.
L'existence d'écoulements à coefficient d'aspect élevé est elle-même liée à des facteurs extérieurs qui renforcent la pertinence d'une description à deux dimensions : la force de gravité induit une forte asymétrie entre l'horizontale et la verticale. On démontre par ailleurs que dans un fluide stationnaire barotrope en rotation (section Deleersnijder), le mouvement est invariant selon l'axe de rotation (\cindex{théorème de Taylor-Proudman}).
Nous utiliserons la notation $\uh$ pour signifier la composante horizontale du mouvement.

\subsection{Écoulements 2-D typiques}
On peut appréhender le monde bi-dimensionnel en se familiarisant avec des écoulements typiques.

\paragraph{Écoulement de translation}
Il s'agit d'un écoulement trivial, $\uh=\vec{c}$.

\paragraph{Écoulement divergent}
On peut modéliser un écoulement globalement divergent par $\uh=c\rh$, où $\rh$ est la composante horizontale du vecteur position, depuis l'origine.
On vérifie (exercice) que $\Div\uh$ vaut dans ce cas $2c$, le facteur~2 étant associé aux deux dimensions de l'écoulement. Cet écoulement implique une divergence sur tout le champ et des vitesses croissant indéfiniment au fur et à mesure que l'on s'éloigne du centre. On peut localiser la structure divergente et lui donner une dimension typique $L$ en imposant : $\uh=ce^{-r_h/L}\rh$. On vérifie aussi que ces écoulements sont irrotationnels.

\paragraph{Rotation solide\label{para:Vortex}}
Par analogie à l'écoulement divergent, on modélise un \cindex{écoulement rotationnel} par $\uh=\vec{\Omega} \cross \rh$.

En coordonnées polaires, cet écoulement s'écrit simplement:
\begin{displaymath}
  u_\theta = \Omega r
\end{displaymath}

Dans l'écoulement 2-D, le vecteur $\vec\Omega$ qui induit la rotation est orthogonal au plan de l'écoulement. On retrouve le vecteur instantané de rotation (la \cindex{vorticité}) en appliquant l'opérateur rotationnel : $\vec\vort \eqdef \curl\uh=2\vec\Omega$.

On note la composante verticale de vorticité
\begin{displaymath}
  \vort_z = 2 \Omega
\end{displaymath}

Ce modèle de rotation solide implique des vitesses croissant indéfiniment à mesure que l'on s'éloigne de l'origine.

\paragraph{Vortex irrotationnel}
Il est possible de créer un flux circulaire qui est non-rotationnel presque partout!

Considérons le mouvement dans le plan, en coordonnées polaires:
\begin{displaymath}
  u_\theta = \frac {\Gamma}{2\pi r}
\end{displaymath}

Cette fois, la vitesse angulaire de rotation n'est plus constante : elle diminue selon une loi de proportionalité inverse.

Pour calculer le rotationnel, on peut utiliser son expression en coordonnées polaires (seule la composante verticale est non-nulle, puisqu'il s'agit d'un mouvement dans le plan):
\begin{displaymath}
  \zeta_z =%
  \frac{1}{r}%
  \left(\pd{(r u_\theta)}{r} - \pd{u_r}{r}\right)
\end{displaymath}

Si on applique le théorème de la circulation:
\begin{displaymath}
  \oint u_\theta\,r\dif \theta = \Gamma = \iint \vort_z \dif A
\end{displaymath}

La vorticité est donc nulle partout, sauf au centre où elle se comporte comme un pic de Dirac.

Bien que ce type de vortex semble tiré par les cheveux, il est en fait très courant car il permet de créer une circulation tout en étant irrotationnel presque partout.

\paragraph{Vortex exponentiel}
Si l'on applique un facteur exponentiel décroissant, on génère une rotation différentielle typique d'un vortex.
\begin{displaymath}
  \vec{u}_h = e^{-r_h/L} \vec{\Omega}\cross\rh
\end{displaymath}
On vérifie que cet écoulement est non-divergent.

\paragraph{Cisaillement}
Le cisaillement $\vec{u}_h = (0,Lx)$ constitue un autre exemple d'écoulement rotationnel, non-divergent (plus généralement, toute forme $(g(y), f(x))$ sera non-divergente).

\subsection{Potentiels scalaires et vecteurs}
Reprenons la décomposition de Helmholtz, mais adaptée au cas bi-dimensionnel. Le potentiel scalaire $\vpot$ est défini dans le plan et on ne prend que son gradient horizontal. Le potentiel vectoriel $\vec{\sfn}$ est un principe à trois dimensions. On peut manipuler l'opérateur rotationnel pour extraire trois composantes :
\begin{equation}
  \begin{split}
     \curl \vec{\sfn} %
                     &= \curl ( \sfn_x \ex  + \sfn_y \ey + \sfn_z \ez ) \\
                     &= \grad \sfn_x \cross \ex +  \grad \sfn_y \cross \ey +  \grad \sfn_z \cross \ez
   \end{split}
\end{equation}
Chaque composante produit un champ vectoriel dans les trois plans orthogonaux aux vecteurs de base. Si on considère un écoulement horizontal, dans le plan $(\ex, \ey)$, alors il ne faut retenir que le troisième terme.

\subsection{Écoulement bi-dimensionnel stationnaire incompressible}
\subsubsection{Fonction de courant \label{sect:courant}}
Si le flux est non-divergent $\Div\u_h = 0$ soit, en développant
\begin{displaymath}
  \pd{u}{x} + \pd{u}{y} = 0,
\end{displaymath}
où l'on a posé $\uh=u\ex + v\ey$. Cette équation est automatiquement rencontrée si $\u_h$ dérive d'un \cindex{potentiel vecteur}. Soit: $\uh = \curl (\sfn_z \ez) = \grad \sfn_z \times \ez$, on a:
\begin{displaymath}
  u = \pd{\sfn_z}{y}\quad \text{et} v = -\pd{\sfn_z}{x}.
\end{displaymath}
L'équation de continuité est satisfaite par $\pd[2]{\sfn_z}{xy} = \pd[2]{\sfn_z}{yx}$.

La fonction $\sfn_z$ porte le nom de \cindex{fonction de courant}, en anglais \emph{streamfunction}, pour les deux raisons suivantes :

\begin{key}{Le flux entre deux points est donné par la différence de la fonction de courant entre ces deux points}
  Sous l'hypothèse d'une masse volumique $\rho$ constante, le flux à travers une section $\mathcal{S}$ quelconque entre deux points vaut
  \begin{displaymath}
    \begin{split}
      \int_a^b \rho \uh\cdot \vec{n}\, \dif {\mathcal S} &= \rho \int_a^b \uh\cdot(\dif y, -\dif x) \\
                                                         &= \int_a^b \pd{\sfn_z}{y}\dif y + \pd{\sfn_z}{x}\dif x  \\
                                                         &= \sfn_{z}(b) - \sfn_{z}(a).
    \end{split}
  \end{displaymath}
\end{key}

Par ailleurs,
\begin{key}{Dans un fluide stationnaire, une particule fluide conserve $\sfn_z$ le long de son mouvement.}
En effet, puisque $\pdt{\psi_z}=0$ (hypothèse stationnaire),
\begin{displaymath}
  \Dif \psi_z = \pdx {\psi_z} \Dif x + \pdy {\psi_z} \Dif y.
\end{displaymath}
\end{key}

Ainsi, dans un écoulement bi-dimensionnel incompressible, les iso-contours de la fonction de courant sont des lignes de courant. Dans ces conditions, deux iso-contours bordent et "piègent" le flux qui ne peut jamais les croiser. La différence entre les valeurs prises par la fonction de courant indique (à un facteur $\rho$ près) le flux de masse entre les deux iso-contours.

\section{Description Lagrangienne d'un écoulement non-stationnaire}
Nous devons introduire deux nouvelles notions, qui traduisent toutes deux la \cindex{description lagrangienne} du champ de vitesse.

\subsection{Ligne de chemin}
La \cindex{ligne de chemin} (\emph{pathline}) est la courbe paramétrique décrivant la trajectoire d’une particule fluide dans un champ de vitesse, possiblement non stationnaire. Elle est définie par :
\begin{equation}
  \od{\vec{p}}{t} = \vec{u}(\vec{p},t), \qquad \vec{p}(t_0) = \vec{p}_0.
  \label{eq:ligne_chemin}
\end{equation}
On peut également l’exprimer à l’aide de la \cindex{fonction de flot}. Dans sa version non-stationnaire, la fonction flot est indexée par deux paramètres, le temps initial ($t_0$) et le temps cible ($t$), et exprime l'advection d'un domaine pris au temps $t_0$, tel qu'il apparaît au temps $t$.
\begin{equation}
  \vec{p}(t) = \phi^{t_0, t }(\vec{x}_0),
  \label{eq:ligne_chemin_flot}
\end{equation}
Le paramètre de la courbe est $t$.
Chaque ligne de chemin correspond donc à la trajectoire d’une particule donnée, définie par sa position initiale $\vec{x}_0$ et son temps initial $t_0$.
À chaque instant $t$, la tangente à la ligne de chemin coïncide avec la direction du champ de vitesse $\vec{u}(\vec{x},t)$ en ce point.
Dans un écoulement laminaire, deux particules lâchées au même instant $t_0$ ne peuvent jamais se collisionner : une telle collision impliquerait qu’en un même point et à un même instant, deux vitesses différentes sont possibles, ce qui contredirait l’hypothèse de continuité du milieu. En revanche, deux lignes de chemins associées à des temps initiaux différents peuvent très bien se croiser.
La ligne de chemin peut donc être interprétée comme une courbe qui se construit dynamiquement au fur et à mesure de l’évolution du temps.

\subsection{Ligne de trace (streakline)}
La \cindex{ligne de trace} (\emph{streakline}) est définie à un instant donné $t$ comme l’ensemble des positions atteintes à cet instant par des particules lâchées à différents instants $t_0$, mais depuis une même position initiale $\vec{x}_0$.
Autrement dit, le paramètre qui permet de se déplacer le long de la courbe est le temps de lâcher $t_0$. En reprenant le formalisme du flot, on écrit :
\begin{equation}
  \vec{s}(t_0) = \phi^{t_0, t}(\vec{x}_0),
  \label{eq:streakline_flot}
\end{equation}
Le paramètre de la courbe est cette fois $t_0$, et non $t$.
En parcourant la courbe, on explore les positions actuelles des particules émises successivement depuis la source $\vec{x}_0$ à différents instants passés.

Ainsi, dans un écoulement non stationnaire, les lignes de courant, les lignes de chemin et les lignes de trace ne coïncident pas nécessairement : chacune décrit un aspect distinct du mouvement du fluide.

\begin{figure}
  \includegraphics[width=\textwidth]{Python/01_vortex_streamlines.pdf}
  \caption{Illustration des lignes de courant (traits pleins), trajectoire (bleue) et ligne de trace (rouge) pour un vortex se déplaçant de gauche à droite. Code interactif sur \url{https://mybinder.org/v2/gh/mcrucifix/lphys1213/main?labpath=IPynb/01_vortex_streamlines.ipynb}} \label{fig:vortex_stream}
\end{figure}

